\documentclass[11pt]{article}
\usepackage{fontspec,amsmath,amssymb,amsthm,hyperref}
\usepackage[margin=1in]{geometry}
\setlength{\emergencystretch}{3em}
\title{Zero-dimensional state spaces have zero transport cost}
\author{LangevinLean mathematical exposition}
\date{}
\begin{document}\maketitle
\noindent Source module: \texttt{UnadjustedLangevin/Degenerate.lean}. Compile this proof document with XeLaTeX or LuaLaTeX. The accompanying HTML pairs the mathematical argument with the exact Lean statements and records the scope of the formalization.


On a subsingleton state space every two states coincide, so every coupling has zero quadratic difference cost. Probability laws on such a space therefore have zero Wasserstein distance.


\section{couplingCost eq zero of subsingleton}


In a subsingleton state space every pair satisfies \(x=y\), so the displacement function is identically zero. Its extended \(L^2\) norm is zero under every pair measure, hence every coupling cost vanishes.



\noindent Lean declarations: \texttt{couplingCost\_eq\_zero\_of\_subsingleton}.

\section{wasserstein2 eq zero of subsingleton}


The independent product of two probabilities is a coupling, so the coupling type is nonempty. Every coupling has zero cost by the preceding result; taking their infimum gives \(W_2(\mu,\pi)=0\).



\noindent Lean declarations: \texttt{wasserstein2\_eq\_zero\_of\_subsingleton}.

\section{wasserstein2Real eq zero of subsingleton}


Rewrite the extended Wasserstein distance as zero using the subsingleton theorem. Its real conversion is also zero. This covers dimension-zero spaces without invoking the nontrivial-space comparison \(m\le L\).



\noindent Lean declarations: \texttt{wasserstein2Real\_eq\_zero\_of\_subsingleton}.

\end{document}
